\documentclass[10pt,a4paper]{amsart}
\usepackage[margin=19mm]{geometry}
\usepackage{amsmath,amssymb,mathtools}
\usepackage[scr=rsfs,cal=euler]{mathalfa}
\usepackage{xfrac}
\usepackage[unicode,hidelinks,bookmarks=false]{hyperref}
\newcommand{\ie}{i.e.\ }
\newcommand{\cf}{cf.\ }
\newcommand{\resp}{respectively\ }
\newcommand{\Subsection}{\subsection}
\providecommand{\nobreakdash}{\nobreak\hbox{-}}
\setlength{\emergencystretch}{2em}
\multlinegap0pt
\usepackage{bm}
\usepackage{bbm}
\usepackage{dsfont}
\usepackage{xspace}
\usepackage{cancel}
\usepackage{mathtools}
\usepackage{amsthm}
\makeatletter
\@ifundefined{theorem}{\newtheorem{theorem}{Theorem}}{}
\@ifundefined{lemma}{\newtheorem{lemma}[theorem]{Lemma}}{}
\makeatother
\renewcommand{\theta}{\vartheta}
\title[$q$-Deformed Topological Recursion, Weight Vectors and Algeb]{$q$-Deformed Topological Recursion, Weight Vectors and Algebraic Structures}
\author{Fridolin Melong, Raimar Wulkenhaar}
\date{Source: arXiv:2606.16441v1 (CC BY 4.0)}
\begin{document}
\maketitle

\noindent\emph{Source and licence.} $q$-Deformed Topological Recursion, Weight Vectors and Algebraic Structures. Version arXiv:2606.16441v1 (CC BY 4.0), distributed under the Creative Commons Attribution 4.0 International licence, as recorded in the arXiv licence metadata for this exact version and shown on the abstract page https://arxiv.org/abs/2606.16441v1. Full rights notice: licenses/arxiv-2606.16441-CC-BY-4.0.txt.\par\smallskip

\noindent\textit{Editorial scope.} Counted unit: the Lemma whose source label is \texttt{l:Aq4\_q} in the article (source0.tex lines 1855--1860, source label \texttt{l:Aq4\_q}), delivered together with the article's own complete original proof of it (source0.tex lines 1861--1901, one \texttt{proof} environment of 41 lines). No mathematics is written, repaired or re-derived: statement and proof are the article's own text line for line. Required context retained uncounted: none. Every definition, display and label of those lines is retained unchanged. Omitted from this package and not redistributed: 1--1854, 1902--2043. Nothing inside a retained excerpt is omitted or altered: every retained body is delivered exactly as the article's lines are written, so no omission receipt of any kind accompanies this package. Citations: the retained excerpts carry no citation command at all, so no bibliography entry of the article is transcribed and no reference is redistributed. Reference adaptation: none was needed and none was made; every reference inside the delivered unit resolves inside it, so no unresolved reference or citation is produced. Printed slips kept literal: no printed word, symbol, hypothesis or number of the counted statement or of its proof was altered. Layout and class: the article's own class is \texttt{amsart} with packages including nicematrix, inputenc, latexsym, pgf, tikz, hyperref, cite, color, mathrsfs, amstext, amsmath, amssymb, amsfonts, bm; the wrapper is the lane's pinned \texttt{amsart} editorial wrapper at 10pt on A4, which restates the theorem-like environments the retained text needs and adds the article's own macro definitions that the retained text needs (\texttt{\textbackslash theta}), with extra wrapper packages (bm, bbm, dsfont, xspace, cancel, mathtools). The delivered numbering is the unit's own. Assets: no retained span contains a figure environment, a graphics-inclusion command, a PDF-page inclusion, a table or a TikZ picture; no image file is copied into this package, the package declares no asset dependency and no image file is redistributed with it. Licence: version arXiv:2606.16441v1 of this article is distributed under the Creative Commons Attribution 4.0 International licence, as recorded in the arXiv licence metadata for this exact version and shown on the abstract page https://arxiv.org/abs/2606.16441v1. A full rights notice is delivered at licenses/arxiv-2606.16441-CC-BY-4.0.txt. Characters: every retained excerpt is pure ASCII, so the delivered engine, XeLaTeX with its default Latin Modern text font, reports no missing character.
				\begin{lemma}[General $q$-deformed kernel evaluation]\label{l:Aq4_q}
					{Let $r \geq 2$, $s > 0$, and $d = \gcd(r, s)$ with $r' = r/d$. The $q$-deformed kernel satisfies the following identity:
						\begin{equation}
							i \Psi^{(0)}_{r,q}\big(-s, \dots, -s, (i-1)s\big) = \big(-1\big)^{i-1} r \big(-1\big)^{\lfloor \frac{i-1}{r'} \rfloor} \binom{d-1}{\lfloor \frac{i-1}{r'} \rfloor}.
					\end{equation}}
				\end{lemma}

				\begin{proof}
					{Let $x_m = \theta^m q^{-2m}$ be the $q$-deformed  variables. We define the sum $\psi_{i,s}^q := i \Psi^{(0)}_{r,q}(-s, \dots, -s, (i-1)s)$. By the definition of the $q$-deformed kernel, this quantity can be expressed in terms of elementary symmetric polynomials:
						\begin{equation}
							\psi_{i,s}^q = \sum_{m=0}^{r-1} x_m^{-(i-1)s} \left( \sum_{\substack{L \subseteq \{0, \dots, r-1\} \setminus \{m\} \\ |L|=i-1}} \prod_{l \in L} x_l^s \right) = \sum_{m=0}^{r-1} x_m^{-(i-1)s} e_{i-1}(\vec{{x}}^s[m]),
						\end{equation}
						where $\vec{{x}}^s = (x_0^s, x_1^s, \dots, x_{r-1}^s)$ represents the $s$-th powers of the $q$-deformed alphabet.}
					
					{The inclusion-exclusion identity for elementary symmetric polynomials $e_i$ states:
						\begin{equation}\label{eq:inc_ex_en}
							e_i(\vec{{x}}^s) = x_m^s e_{i-1}(\vec{{x}}^s[m]) + e_i(\vec{{x}}^s[m]).
						\end{equation}
						Multiplying \eqref{eq:inc_ex_en} by $x_m^{-is}$ and summing over the index $m \in \{0, \dots, r-1\}$ yields:
						\begin{equation}
							e_i(\vec{{x}}^s) \sum_{m=0}^{r-1} x_m^{-is} = \sum_{m=0}^{r-1} x_m^{-(i-1)s} e_{i-1}(\vec{{x}}^s[m]) + \sum_{m=0}^{r-1} x_m^{-is} e_i(\vec{{x}}^s[m]).
						\end{equation}
						Recognizing $\psi_{i,s}^q$ in the first term and $\psi_{i+1,s}^q$ in the second term on the right-hand side, we obtain the recurrence:
						\begin{equation}\label{eq:rec_q_en}
							e_i(\vec{{x}}^s) \left( \sum_{m=0}^{r-1} x_m^{-is} \right) = \psi_{i,s}^q + \psi_{i+1,s}^q.
					\end{equation}}
					
					{The variables are given by $x_m^s = (\theta^s q^{-2s})^m$. Since $d = \gcd(r, s)$, the set $\{x_m^s\}_{m=0}^{r-1}$ consists of the $r'$ roots of the equation $t^{r'} - (q^{-2s})^{r'} = 0$, each occurring with multiplicity $d$. The generating function for $e_i(\vec{{x}}^s)$ is:
						\begin{equation}
							\sum_{i=0}^r e_i(\vec{{x}}^s) t^i = \prod_{a=0}^{r-1} (1 + t x_a^s) = \left( 1 + (-1)^{r'} C_q t^{r'} \right)^d,
						\end{equation}
						where $C_q = \prod_{k=0}^{r'-1} x_k^s$. Expanding this via the binomial theorem shows that $e_i(\vec{{x}}^s) \neq 0$ only when $i$ is a multiple of $r'$. Specifically, for $i = j r'$, $e_{jr'}(\vec{{x}}^s) = \binom{d}{j} (-1)^{j r'} C_q^j$.}
					
					{The sum $S_i = \sum_{m=0}^{r-1} x_m^{-is}$ vanishes unless $is \equiv 0 \pmod{r}$, which corresponds to $i$ being a multiple of $r'$. For $i = j r'$, the term $x_m^{-is}$ compensates for the $q$-dependent factor $C_q^j$ in $e_i(\vec{{x}}^s)$. Consequently, the $q$-dependence drops out of the recurrence \eqref{eq:rec_q_en}, leading to:
						\begin{equation}
							\psi_{i,s}^q + \psi_{i+1,s}^q = r \delta_{r'|i} (-1)^{i/r'} \binom{d}{i/r'}.
					\end{equation}}
					
					{Solving the recurrence with the initial condition $\psi_{1,s}^q = r$ (the spin-1 mode) yields:
						\begin{equation}
							\psi_{i,s}^q = (-1)^{i-1} r \sum_{j=0}^{\lfloor (i-1)/r' \rfloor} (-1)^j \binom{d}{j}.
						\end{equation}
						Applying the binomial identity $\sum_{j=0}^k (-1)^j \binom{d}{j} = (-1)^k \binom{d-1}{k}$, we arrive at:
						\begin{equation}
							\psi_{i,s}^q = (-1)^{i-1} r (-1)^{\lfloor \frac{i-1}{r'} \rfloor} \binom{d-1}{\lfloor \frac{i-1}{r'} \rfloor}.
						\end{equation}
						This completes the proof.}
				\end{proof}

\end{document}
